% Einheit 3 von 3 – Zählen mit Bedingungen und Wahrscheinlichkeitsterme (bank/kombinatorik/e3.jsonl, zusammenbau v0.1)
%% TODO Zeitmarke Einheit 3: Marken-Zeile nicht lesbar
%% TODO Zweigzeile Teil 1: was hier gelernt wird (Satz in Schülersprache aus dem Einheitstitel) – Einheit 3
\einheitenkopf[e3]{Einheit 3 von 3 · Zählen mit Bedingungen und Wahrscheinlichkeitsterme}
\zweigzeile{Abitur GK}
\setcounter{aufgabe}{11}

%% TODO Titel: Ich-kann-Satz fehlt in der Bank (Platzhalter: Kettenname) – Nr. 12
%% TODO Anweisung: ein Satz über der ersten Teilaufgabe fehlt in der Bank – Nr. 12
\begin{aufgabe}{Zählen mit Bedingungen und Wahrscheinlichkeitsterme}
%% TODO \rechenplatz{n}: Schrittzahl des Grundfalls fehlt in der Bank (nur bei fester Schreibform, 2.3 b)
\begin{teile}
\teil Ein Strauß besteht entweder aus zwei roten oder aus drei gelben Tulpen; gezählt werden alle möglichen Sträuße. Werden die Anzahlen der beiden Fälle multipliziert oder addiert? \\ \kreuz{mal} \\ \kreuz{plus}
\teil In einer Reihe stehen fünf Stühle. Zwei Stühle werden so ausgewählt, dass sie nicht nebeneinander stehen. Schreibe alle Auswahlen auf. Begründe, dass keine fehlt. \leerfeld
\teil In einer Reihe stehen fünf Stühle. Ali und Bea setzen sich so, dass sie nicht nebeneinander sitzen. Wie viele Möglichkeiten gibt es? \leerfeld
\teil Ein Programm hat sechs Punkte: zwei Reden und vier Musikstücke, alle verschieden. Die beiden Reden sollen nicht direkt nacheinander kommen. Wie viele Reihenfolgen gibt es \hfill (Abitur 2026 GK)? \leerfeld
\teil Ein Strauß besteht aus elf Rosen in den Farben Rot, Weiß und Gelb; von jeder Farbe sind mindestens drei und höchstens fünf dabei. Wie viele Möglichkeiten gibt es, den Strauß zusammenzustellen \hfill (Abitur 2020 LK)? \leerfeld
\teil Ein Würfel wird viermal geworfen. Wie groß ist die Wahrscheinlichkeit, dass lauter verschiedene Augenzahlen fallen? $P = $ \leerfeld
\teil Ein Glücksrad mit den drei gleich großen Feldern 1, 2 und 3 wird viermal gedreht. Gib einen Term für die Wahrscheinlichkeit an, dass jede Zahl mindestens einmal kommt, und berechne ihn \hfill (Abitur 2023 LK). $P = $ \leerfeld
\teil Ein Tetraeder mit den Zahlen 1 bis 4 wird achtmal geworfen. Gib einen Term für die Wahrscheinlichkeit an, dass nicht jede Zahl genau zweimal fällt, und berechne ihn auf drei Stellen \hfill (Abitur 2022 LK). $P \approx $ \leerfeld
\end{teile}
\end{aufgabe}

%% TODO Titel: Ich-kann-Satz fehlt in der Bank (Platzhalter: Kettenname) – Nr. 13
%% TODO Anweisung: ein Satz über der ersten Teilaufgabe fehlt in der Bank – Nr. 13
\begin{aufgabe}{Zählen mit Bedingungen und Wahrscheinlichkeitsterme – Fehler finden, Begründen, Anwendung}
\begin{teile}
\teil Sieben Plätze liegen in einer Reihe; zwei werden so ausgewählt, dass mindestens drei Plätze dazwischen liegen. Tim zählt so. Finde den Fehler und gib alle Auswahlen an. \rechnung{&1 \text{ und } 5, \ 1 \text{ und } 6, \ 1 \text{ und } 7 \\ &\text{also } 3 \text{ Auswahlen}}
\teil Begründe, warum man bei Sitzordnungen mit einer Abstandsbedingung erst die Muster aufzählt und dann multipliziert.
\teil Ein Fotograf stellt vier Erwachsene und drei Kinder in eine Reihe; keine zwei Kinder sollen nebeneinander stehen. Wie viele Aufstellungen gibt es? \leerfeld
\end{teile}
\end{aufgabe}
